%---
% title: "Ma trận và phép biến đổi tuyến tính"
% date: 2026-08-05
% tags: [toan, dai-so]
% summary: "Ma trận biểu diễn phép biến đổi tuyến tính như thế nào, kèm ví dụ."
%---
\section{Ma trận}

Ma trận là bảng số hình chữ nhật. Mỗi ma trận \(A\) kích thước
\(m \times n\) có \(m\) hàng và \(n\) cột.

\begin{equation}
A = \begin{pmatrix} 1 & 2 \\ 3 & 4 \end{pmatrix}
\end{equation}

\section{Phép biến đổi tuyến tính}

Mỗi ma trận xác định một ánh xạ tuyến tính từ \(\mathbb{R}^n\) vào
\(\mathbb{R}^m\).

\section{Ví dụ}

Ma trận xoay trong mặt phẳng một góc \(\theta\):

\begin{equation}
R(\theta) = \begin{pmatrix} \cos\theta & -\sin\theta \\
\sin\theta & \cos\theta \end{pmatrix}
\end{equation}

\section{Kết luận}

Ma trận là cầu nối giữa đại số tuyến tính và hình học.
